\documentclass[11pt, a4paper]{article}
\usepackage[dutch]{babel}
\usepackage[margin=2cm]{geometry}
\usepackage[european]{circuitikz}
\usepackage{amsmath}
\usepackage{enumitem}

\begin{document}

% --- bladzijde 1: hoofding en strategie ---

\begin{center}
    \Large \textbf{Taak fysica: schakelingen}
\end{center}
\vspace{1cm}

\textbf{Naam:} \hrulefill \hfill \textbf{Klas:} \hrulefill \hfill \textbf{Datum:} \hrulefill 
\vspace{2cm}

\noindent\fbox{
    \begin{minipage}{\textwidth}
    \vspace{0.2cm}
    \textbf{oplossingsstrategie voor schakelingen}
    \begin{enumerate}[leftmargin=*]
        \item \textbf{symbolen noteren:} noteer bij elke component de symbolen $U$, $I$ en $R$. vul de bekende waarden uit de opgave in.
        \item \textbf{de '2 van de 3' regel:} als van een component twee van de drie grootheden gekend zijn, bereken je de derde via de wet van ohm ($U = I \cdot R$).
        \item \textbf{vereenvoudigen:} vervang stukken die zuiver serie of parallel staan door hun substitutieweerstand. herhaal dit tot je de totale weerstand vindt.
        \item \textbf{spanning en stroom verdelen:} 
        \begin{itemize}
            \item \textbf{serie:} de stroom $I$ is overal gelijk. de spanning $U$ verdeelt zich: de grootste spanning staat over de grootste weerstand.
            \item \textbf{parallel:} de spanning $U$ is overal gelijk. de stroom $I$ verdeelt zich: de grootste stroom gaat door de kleinste weerstand.
        \end{itemize}
    \end{enumerate}
    \vspace{0.2cm}
    \end{minipage}
}

\newpage
% --- bladzijde 2: oefening 1 ---

\textbf{oefening 1} \\
een stroomkring bevat een spanningsbron van $12\text{ V}$. de schakeling bestaat uit een weerstand $R_1 = 20\,\Omega$ in serie met een parallelblok van $R_2 = 30\,\Omega$ en $R_3 = 60\,\Omega$.

\begin{enumerate}[label=\alph*)]
    \item teken het elektrisch schema en benoem alle componenten.
    \vspace{7.5cm}
    \item bereken de substitutieweerstand van de volledige schakeling.
    \vspace{7.5cm}
    \item bereken de totale stroomsterkte die de bron levert.
    \vfill
\end{enumerate}

\newpage
% --- bladzijde 3: oefening 2 ---

\textbf{oefening 2} \\
in de onderstaande schakeling levert de bron een spanning $U = 24\text{ V}$. de weerstandswaarden zijn: $R_1 = 10\,\Omega$, $R_2 = 30\,\Omega$ en $R_3 = 20\,\Omega$.

\vspace{1cm}
\begin{center}
    \begin{circuitikz}[scale=1.4]
        \draw (0,0) to[battery1, l=$U{=}24\text{ V}$] (0,3)
              to[short] (2,3)
              to[R, l=$R_1$, a=$10\,\Omega$] (4,3)
              to[R, l=$R_2$, a=$30\,\Omega$] (6,3)
              to[short] (8,3)
              to[short] (8,0)
              to[short] (0,0);
        \draw (2,3) to[short, *-] (2,1.5)
              to[R, l=$R_3$, a=$20\,\Omega$] (6,1.5)
              to[short, -*] (8,1.5);
    \end{circuitikz}
\end{center}
\vspace{1cm}

\begin{enumerate}[label=\alph*)]
    \item bepaal de stroomsterkte door weerstand $R_3$.
    \vspace{8.5cm}
    \item bereken de spanning over weerstand $R_1$ en de spanning over weerstand $R_2$.
    \vfill
\end{enumerate}

\newpage
% --- bladzijde 4: oefening 3 ---

\textbf{oefening 3} \\
gegeven is de onderstaande gemengde schakeling met $R_1 = 20\,\Omega$, $R_2 = 30\,\Omega$, $R_3 = 8\,\Omega$, $R_4 = 60\,\Omega$ en een bronspanning $U = 60\text{ V}$.

\vspace{0.3cm}
\begin{center}
    \begin{circuitikz}[scale=1.1] % schaal iets verkleind naar 1.1
        \draw (0,0) to[battery1, l=$U{=}60\text{ V}$] (0,4)
              to[short] (1,4)
              to[short, *-] (1,5)
              to[R, l=$R_1$, a=$20\,\Omega$] (4,5)
              to[short, -*] (4,4)
              to[R, l=$R_3$, a=$8\,\Omega$] (7,4)
              to[short] (9,4)
              to[short] (9,0)
              to[short] (0,0);
        \draw (1,4) to[short] (1,3)
              to[R, l=$R_2$, a=$30\,\Omega$] (4,3)
              to[short] (4,4);
        \draw (1,4) to[short, *-] (1,1.5)
              to[R, l=$R_4$, a=$60\,\Omega$] (7,1.5)
              to[short, -*] (9,1.5);
    \end{circuitikz}
\end{center}
\vspace{0.3cm}

\begin{enumerate}[label=\alph*)]
    \item bereken de totale substitutieweerstand.
    \vspace{3.5cm} % witruimte aangepast naar 3.5cm
    \item bepaal de totale stroomsterkte die de bron verlaat.
    \vspace{3.5cm}
    \item bereken de stroomsterkte door $R_4$.
    \vspace{3.5cm}
    \item bereken de stroomsterkte door $R_1$ en door $R_2$.
    \vfill
\end{enumerate}

\end{document}